\documentclass{amsart}
\usepackage{amsmath,amssymb,amsthm,hyperref}
\newtheorem{theorem}{Theorem}
\newtheorem{lemma}{Lemma}
\theoremstyle{definition}
\newtheorem{remark}{Remark}
\begin{document}

\title{Approximate fixed points of VMO self-maps of the ball}
\maketitle

\begin{theorem}
Let $N\ge 1$ and let $B=\overline{B(0,1)}\subset\mathbb{R}^N$ be the closed unit
ball, equipped with Lebesgue measure $\mathcal{L}^N$. Let $f:B\to B$ be a VMO map.
Then for every $\varepsilon>0$ the approximate fixed-point set
\[
A_\varepsilon(f):=\{x\in B:\ |f(x)-x|<\varepsilon\}
\]
has strictly positive Lebesgue measure.
\end{theorem}

Throughout, $\omega_N:=\mathcal{L}^N(B(0,1))$ denotes the volume of the unit ball,
$|E|:=\mathcal{L}^N(E)$, and $\fint_E g:=\frac1{|E|}\int_E g$. Recall that
$f:B\to\mathbb{R}^N$ is VMO means that its \emph{VMO modulus}
\begin{equation}\label{eq:vmo}
\omega(r):=\sup_{\substack{x\in B,\ 0<\rho\le r}}
\ \fint_{B(x,\rho)\cap B}\Bigl|f(y)-\fint_{B(x,\rho)\cap B} f\Bigr|\,dy
\end{equation}
satisfies $\omega(r)\to 0$ as $r\to0^+$. The quantity $\omega(r)$ is nondecreasing
in $r$, and it is finite for $r\le 1$ because $|f|\le 1$ on $B$ (as $f$ maps into
$B$), so the integrand in \eqref{eq:vmo} is bounded by $2$.

Throughout, $f:B\to B$ is fixed and $\varepsilon>0$ is given. We argue by
contradiction: \emph{assume}
\begin{equation}\label{eq:assume}
|A_\varepsilon(f)|=0,\qquad\text{equivalently}\qquad
|f(x)-x|\ge\varepsilon\ \text{ for a.e.\ } x\in B.
\end{equation}

\section{A geometric lower bound for ball--cap volumes}

\begin{lemma}\label{lem:cap}
Let $\kappa_N:=\dfrac{|\{z\in B(0,1):\ z_1\le -\tfrac12\}|}{\omega_N}\in(0,1)$.
Then for every $x\in B$ and every $r$ with $0<r\le 1$,
\[
|B(x,r)\cap B|\ \ge\ \kappa_N\,\omega_N\,r^N\ =\ \kappa_N\,|B(x,r)|.
\]
\end{lemma}

\begin{proof}
The number $\kappa_N$ is a strictly positive constant in $(0,1)$, being the
Lebesgue measure of a nonempty open subset of the unit ball divided by $\omega_N$.

\emph{Case 1: $|x|\le 1-r$.} Then $B(x,r)\subseteq B$, so
$|B(x,r)\cap B|=|B(x,r)|=\omega_N r^N\ge\kappa_N\omega_N r^N$ since $\kappa_N\le1$.

\emph{Case 2: $1-r<|x|\le 1$.} In particular $x\ne0$; set $e:=x/|x|$, a unit
vector. Define the slab
\[
E:=\Bigl\{y\in B(x,r):\ \langle y-x,\,e\rangle\le -\tfrac{r}{2}\Bigr\}.
\]
We first show $E\subseteq B$. For $y\in E$ we have $|y-x|\le r$ and
$\langle y-x,e\rangle\le -r/2$, hence, using $\langle y-x,x\rangle=|x|\,\langle
y-x,e\rangle$,
\[
|y|^2=|x|^2+2\langle y-x,x\rangle+|y-x|^2
\le |x|^2+2|x|\bigl(-\tfrac r2\bigr)+r^2
=|x|^2-|x|r+r^2 .
\]
Consider $g(s):=s^2-sr+r^2$ for $s\in[0,1]$. Being a convex function, $g$ attains
its maximum on $[0,1]$ at an endpoint; since $g(0)=r^2\le1$ and
$g(1)=1-r+r^2=1-r(1-r)\le1$ (because $0\le r(1-r)$ for $r\in[0,1]$), we get
$g(s)\le1$ for all $s\in[0,1]$. With $s=|x|\le1$ this yields $|y|^2\le g(|x|)\le1$,
so $y\in B$. Thus $E\subseteq B(x,r)\cap B$.

Finally we compute $|E|$. Writing $y=x+rz$ with $z\in B(0,1)$, the condition
$\langle y-x,e\rangle\le -r/2$ becomes $\langle z,e\rangle\le -1/2$. By rotational
invariance of Lebesgue measure on the unit ball,
\[
|E|=r^N\,\bigl|\{z\in B(0,1):\ \langle z,e\rangle\le -\tfrac12\}\bigr|
   =r^N\,\bigl|\{z\in B(0,1):\ z_1\le -\tfrac12\}\bigr|
   =\kappa_N\,\omega_N\,r^N .
\]
Therefore $|B(x,r)\cap B|\ge|E|=\kappa_N\omega_N r^N$, completing Case 2.
\end{proof}

\section{The averaging operator}

For $0<r\le1$ define
\begin{equation}\label{eq:fr}
f_r(x):=\fint_{B(x,r)\cap B} f
       =\frac{1}{|B(x,r)\cap B|}\int_{B(x,r)\cap B}f(y)\,dy,\qquad x\in B.
\end{equation}
By Lemma~\ref{lem:cap} the denominator is bounded below by $\kappa_N\omega_N r^N>0$,
so $f_r$ is well defined on $B$.

\begin{lemma}\label{lem:basic}
Fix $0<r\le1$. Then:
\begin{enumerate}
\item[(a)] $f_r:B\to B$ (i.e.\ $f_r$ maps $B$ into $B$);
\item[(b)] $f_r$ is continuous on $B$.
\end{enumerate}
\end{lemma}

\begin{proof}
(a) For each $x$, $f_r(x)$ is the average of $f$ over $B(x,r)\cap B$. Since
$f(y)\in B$ for all $y$ and $B$ is closed and convex, the average of $f$ over any
set of positive finite measure lies in $B$ (averages of values in a closed convex
set belong to that set). Hence $f_r(x)\in B$.

(b) Write $V(x):=|B(x,r)\cap B|$ and $I(x):=\int_{B(x,r)\cap B}f(y)\,dy
=\int_B \mathbf 1_{B(x,r)}(y)\,f(y)\,dy$, so $f_r=I/V$. The map
$x\mapsto \mathbf 1_{B(x,r)}$ is continuous from $B$ into $L^1(B)$: indeed, for
$x\to x'$,
\[
\int_B|\mathbf 1_{B(x,r)}(y)-\mathbf 1_{B(x',r)}(y)|\,dy
=|B(x,r)\,\triangle\,B(x',r)|\to0,
\]
because the boundary sphere $\partial B(x',r)$ has Lebesgue measure zero and
$\mathbf 1_{B(x,r)}\to\mathbf 1_{B(x',r)}$ pointwise off $\partial B(x',r)$, so by
dominated convergence the symmetric-difference volume tends to $0$. Since
$|f|\le1$, $|I(x)-I(x')|\le |B(x,r)\triangle B(x',r)|\to0$, so $I$ is continuous;
taking $f\equiv1$ shows $V$ is continuous. As $V\ge\kappa_N\omega_N r^N>0$ on $B$,
$f_r=I/V$ is continuous.
\end{proof}

\begin{lemma}[Oscillation bound]\label{lem:osc}
Let $C_N:=2^{N}/\kappa_N$. Fix $r$ with $0<r\le\tfrac12$. Let $x_0\in B$ and put
$c:=\fint_{B(x_0,2r)\cap B}f$. Then for every $y\in B(x_0,r)\cap B$,
\begin{equation}\label{eq:osc}
|f_r(y)-c|\ \le\ C_N\,\omega(2r).
\end{equation}
Consequently $|f_r(y)-f_r(y')|\le 2C_N\,\omega(2r)$ for all $y,y'\in
B(x_0,r)\cap B$.
\end{lemma}

\begin{proof}
Let $y\in B(x_0,r)\cap B$. If $z\in B(y,r)$ then
$|z-x_0|\le|z-y|+|y-x_0|<r+r=2r$, so
\begin{equation}\label{eq:incl}
B(y,r)\cap B\ \subseteq\ B(x_0,2r)\cap B.
\end{equation}
Using \eqref{eq:fr} and that $c$ is a constant,
\[
|f_r(y)-c|=\Bigl|\fint_{B(y,r)\cap B}(f-c)\Bigr|
\le\fint_{B(y,r)\cap B}|f-c|
=\frac{1}{|B(y,r)\cap B|}\int_{B(y,r)\cap B}|f-c|.
\]
By \eqref{eq:incl} the last integral is at most
$\int_{B(x_0,2r)\cap B}|f-c|$. Since $c=\fint_{B(x_0,2r)\cap B}f$ and $2r\le1$,
the definition \eqref{eq:vmo} (with the admissible choice $x=x_0$, $\rho=2r$)
gives
\[
\int_{B(x_0,2r)\cap B}|f-c|
=|B(x_0,2r)\cap B|\,\fint_{B(x_0,2r)\cap B}\Bigl|f-\fint_{B(x_0,2r)\cap B}f\Bigr|
\le |B(x_0,2r)\cap B|\,\omega(2r).
\]
Now $|B(x_0,2r)\cap B|\le|B(x_0,2r)|=\omega_N(2r)^N=2^N\omega_N r^N$, while by
Lemma~\ref{lem:cap} $|B(y,r)\cap B|\ge\kappa_N\omega_N r^N$. Combining,
\[
|f_r(y)-c|\le\frac{2^N\omega_N r^N}{\kappa_N\omega_N r^N}\,\omega(2r)
=\frac{2^N}{\kappa_N}\,\omega(2r)=C_N\,\omega(2r),
\]
which is \eqref{eq:osc}. The final assertion follows by applying \eqref{eq:osc} to
both $y$ and $y'$ and using $|f_r(y)-f_r(y')|\le|f_r(y)-c|+|c-f_r(y')|$.
\end{proof}

\section{Brouwer fixed point of $f_r$ and a local $L^1$ estimate}

\begin{lemma}\label{lem:brouwer}
For each $0<r\le1$ there exists $x_r\in B$ with $f_r(x_r)=x_r$.
\end{lemma}

\begin{proof}
By Lemma~\ref{lem:basic}, $f_r:B\to B$ is continuous and $B$ is compact and convex
(it is the closed unit ball of $\mathbb{R}^N$). Brouwer's fixed point theorem
yields $x_r\in B$ with $f_r(x_r)=x_r$.
\end{proof}

\begin{lemma}[Local $L^1$ closeness of $f$ and $f_r$]\label{lem:loc}
Let $0<r\le\tfrac12$ and let $x_0\in B$; set $D:=B(x_0,r)\cap B$. Then
\[
\int_{D}|f-f_r|\ \le\ C'_N\,r^N\,\omega(2r),
\qquad C'_N:=\omega_N\,2^N\Bigl(1+\frac1{\kappa_N}\Bigr).
\]
\end{lemma}

\begin{proof}
Set $c:=\fint_{B(x_0,2r)\cap B}f$. By the triangle inequality,
\[
\int_{D}|f-f_r|\le\int_{D}|f-c|+\int_{D}|f_r-c|.
\]
Since $D=B(x_0,r)\cap B\subseteq B(x_0,2r)\cap B$, the first integral satisfies, by
the same computation as in the proof of Lemma~\ref{lem:osc},
\[
\int_{D}|f-c|\le\int_{B(x_0,2r)\cap B}|f-c|\le |B(x_0,2r)\cap B|\,\omega(2r)
\le 2^N\omega_N r^N\,\omega(2r).
\]
For the second integral, Lemma~\ref{lem:osc} gives $|f_r-c|\le C_N\omega(2r)$
pointwise on $D$, while $|D|\le|B(x_0,r)|=\omega_N r^N$; hence
\[
\int_{D}|f_r-c|\le |D|\,C_N\,\omega(2r)\le \omega_N r^N\cdot\frac{2^N}{\kappa_N}\,\omega(2r).
\]
Adding the two bounds,
\[
\int_{D}|f-f_r|\le\Bigl(2^N\omega_N+\frac{2^N\omega_N}{\kappa_N}\Bigr)r^N\omega(2r)
=\omega_N 2^N\Bigl(1+\frac1{\kappa_N}\Bigr)r^N\omega(2r)=C'_N\,r^N\,\omega(2r). \qedhere
\]
\end{proof}

\section{Conclusion: deriving the contradiction}

Since $\omega(2r)\to0$ and $r\to0$ as $r\to0^+$, we may choose $r_0\in(0,\tfrac12]$
so small that
\begin{equation}\label{eq:choice}
2\,C_N\,\omega(2r)+r\ \le\ \frac{\varepsilon}{2}\qquad\text{for all }0<r\le r_0 .
\end{equation}

Fix any $r$ with $0<r\le r_0$ and let $x_r$ be the Brouwer fixed point from
Lemma~\ref{lem:brouwer}, so $f_r(x_r)=x_r$. Put $D_r:=B(x_r,r)\cap B$. For
$y\in D_r$ we have $|y-x_r|\le r$, so by the last assertion of Lemma~\ref{lem:osc}
(applied with $x_0=x_r$, and the pair $y,x_r\in D_r$),
\[
|f_r(y)-y|\le|f_r(y)-f_r(x_r)|+|f_r(x_r)-y|
=|f_r(y)-f_r(x_r)|+|x_r-y|
\le 2C_N\,\omega(2r)+r .
\]
Combining with \eqref{eq:choice},
\begin{equation}\label{eq:near}
|f_r(y)-y|\ \le\ \frac{\varepsilon}{2}\qquad\text{for all }y\in D_r .
\end{equation}

By \eqref{eq:assume}, for a.e.\ $y\in B$ (in particular for a.e.\ $y\in D_r$) we
have $|f(y)-y|\ge\varepsilon$. Hence, for a.e.\ $y\in D_r$, by \eqref{eq:near},
\[
|f(y)-f_r(y)|\ \ge\ |f(y)-y|-|f_r(y)-y|\ \ge\ \varepsilon-\frac{\varepsilon}{2}
=\frac{\varepsilon}{2} .
\]
Integrating over $D_r$ and using Lemma~\ref{lem:cap} (with $x=x_r$) for the volume,
\begin{equation}\label{eq:lhs}
\int_{D_r}|f-f_r|\ \ge\ \frac{\varepsilon}{2}\,|D_r|
\ \ge\ \frac{\varepsilon}{2}\,\kappa_N\,\omega_N\,r^N .
\end{equation}
On the other hand, Lemma~\ref{lem:loc} (with $x_0=x_r$) gives
\begin{equation}\label{eq:rhs}
\int_{D_r}|f-f_r|\ \le\ C'_N\,r^N\,\omega(2r).
\end{equation}
Comparing \eqref{eq:lhs} and \eqref{eq:rhs} and dividing through by $r^N>0$,
\[
\frac{\varepsilon}{2}\,\kappa_N\,\omega_N\ \le\ C'_N\,\omega(2r),
\qquad\text{i.e.}\qquad
\omega(2r)\ \ge\ \frac{\varepsilon\,\kappa_N\,\omega_N}{2\,C'_N}\ =:\ \delta(\varepsilon).
\]
Because $\kappa_N,\omega_N,C'_N>0$ and $\varepsilon>0$, the constant
$\delta(\varepsilon)$ is strictly positive and \emph{independent of $r$}. Thus
\[
\omega(2r)\ \ge\ \delta(\varepsilon)>0\qquad\text{for all }0<r\le r_0 ,
\]
which, on letting $r\to0^+$, contradicts the VMO hypothesis $\omega(2r)\to0$.

Therefore the assumption \eqref{eq:assume} is untenable, i.e.\
$|A_\varepsilon(f)|>0$. As $\varepsilon>0$ was arbitrary, the theorem is proved.
\qed

\begin{remark}
No quantitative decay rate on $\omega$ is needed. The contradiction is structural:
the lower bound \eqref{eq:lhs} and the upper bound \eqref{eq:rhs} for
$\int_{D_r}|f-f_r|$ scale identically in $r$ (both proportional to $r^N$), so the
powers of $r$ cancel and one is left with a strictly positive constant bounded
above by $\omega(2r)$, which is impossible as $r\to0$.
\end{remark}

\begin{remark}
The heart of the matter is Lemma~\ref{lem:osc}: a VMO map is ``almost continuous''
in the precise sense that its averages $f_r$ have oscillation at most
$2C_N\,\omega(2r)$ on balls of the \emph{same} radius $r$ used to define $f_r$.
Crucially this bound contains \emph{no} factor $1/r$ (which the naive Lipschitz
estimate from $|f|\le1$ would produce). Hence the Brouwer fixed point $x_r$ of
$f_r$ pins $f_r$ near the identity not just at one point but on the whole ball
$D_r$ of volume $\gtrsim r^N$; since $f$ and $f_r$ are $L^1$-close on $D_r$ at the
matching scale, $f$ itself must be within $\varepsilon$ of the identity on a
positive-measure subset of $D_r$. For continuous $f$ (which is VMO) Brouwer already
gives an exact fixed point and a neighborhood of approximate ones; the content here
is the extension to the genuinely discontinuous VMO class, in which jump
discontinuities are precisely what the VMO condition forbids.
\end{remark}

\end{document}
